\chapter*{Resumo}
\pagenumbering{gobble}

O estudo deste Trabalho de Conclusão de Curso tem como objetivo analisar os códigos produzidos por Grandes Modelos de Linguagem na resolução de problemas de nível sênior na Olimpíada Brasileira de Informática, Modalidade Programação. Devido a isso, foram feitos experimentos com cinco linguagens diferentes, utilização de imagens e descrições, informando restrições de tempo e memória e, por fim, técnicas de \textit{prompt}.
Para isso, criou-se uma ferramenta desenvolvida em Python com o objetivo de automatizar a geração de código através de requisições via API dos modelos, verificando se as respostas produzidas pelos modelos contêm um código executável para o \textit{Judge} compilar e executar no próprio ambiente da máquina.
Além disso, todas as requisições para os LLMs foram feitas com temperatura igual a zero, mantendo a característica de decodificação gulosa (\textit{greedy decoding}), ou seja, diminuir a variação dos códigos produzidos pelos modelos. Isso impacta diretamente na escolha da métrica $\passat{1}$ ($\passat{k}$, tal que $k = 1$) para os vereditos do \textit{Judge}, pois o objetivo do estudo não é avaliar a ``criatividade'' dos modelos na resolução de problemas de nível sênior, mas sim a capacidade de resolução dos mesmos problemas em cenários diferentes.
Assim, o estudo mostra que, nesse \textit{dataset}, a complexidade sintática das linguagens tem impacto direto na resolução do problema; além disso, mesmo utilizando a melhor linguagem (C++), que melhor performou na média dos modelos, o uso de imagens no formato Base64 nos \textit{prompts} teve impacto direto no consumo de \textit{tokens} e não garante a resolução do problema. Entretanto, na linguagem que teve o pior desempenho (C), mesmo informando o tempo limite de execução e a memória limite, as soluções produzidas não melhoraram as respostas reproduzidas pelos LLMs. Além disso, a utilização de técnicas de \textit{prompt} ajudou os modelos a melhorarem seu desempenho no $\passat{1}$ ao informar alguns exemplos simples de entrada e saída. Por fim, mostra-se que os grandes modelos de linguagem ainda possuem grandes desafios para a produção de soluções de problemas de programação, pois a lógica para resolver a questão é a mesma, porém a predição do próximo \textit{token} tem impacto direto, principalmente em linguagens que precisam de implementações de estruturas de dados, otimizações, etc.

\vspace{5mm}
\noindent\textbf{\textit{Palavras-chave}:~}Modelos de Linguagem de Grande Porte; Programação Competitiva; Olimpíada Brasileira de Informática; Avaliação Experimental; Engenharia de \textit{Prompt}; Visão Multimodal.
